\chapter{Evaluation und Ergebnisse}
\label{chap:evaluation_und_ergebnisse}

\section{Leistungs-Analyse der Transformer- und LightGBM-Modelle}

Tabelle \ref{tab:overfitting_check} vergleicht die Accuracy- und F1-Werte der Trainings- und Testphase für alle fünf untersuchten Transformer-Modelle. 
Die zusätzlich angegebenen Differenzwerte ($\Delta$) zeigen den Unterschied zwischen Test- und Trainingsergebnissen. 
Da die Abweichungen in beiden Metriken bei allen Modellen sehr klein ist, lässt sich daraus schließen, dass kein Overfitting vorliegt. 
Die Modelle generalisieren gut und zeigen eine stabile Leistung auf bislang ungesehenen Daten.

\begin{table}[ht]
\centering
\begin{tabular}{lcccccc}
    \toprule
    \multirow{2}{*}{Modell} & 
    \multicolumn{2}{c}{Training} & 
    \multicolumn{2}{c}{Test} & 
    \multicolumn{2}{c}{$\Delta$ (Test - Train)} \\
    \cmidrule(lr){2-3} \cmidrule(lr){4-5} \cmidrule(l){6-7}
    & Accuracy & F1 & Accuracy & F1 & $\Delta$ Acc & $\Delta$ F1 \\
    \midrule
    XLM-RoBERTa-Large   & 0.9780 & 0.9780 & 0.9795 & 0.9795 & +0.0015 & +0.0015 \\
    XLM-RoBERTa-Base    & 0.9730 & 0.9729 & 0.9717 & 0.9716 & $-$0.0013 & $-$0.0013 \\
    RoBERTa-Large       & 0.9795 & 0.9795 & 0.9765 & 0.9764 & $-$0.0030 & $-$0.0031 \\
    RoBERTa-Base        & 0.9771 & 0.9771 & 0.9751 & 0.9751 & $-$0.0020 & $-$0.0020 \\
    BERT-Base-Uncased   & 0.9507 & 0.9505 & 0.9533 & 0.9531 & +0.0026 & +0.0026 \\
    \bottomrule
\end{tabular}
\caption{Vergleich von Training und Test: Accuracy und F1 zur Überprüfung von Overfitting}
\label{tab:overfitting_check}
\end{table}

Tabelle \ref{tab:vergleich_der_transformer_modelle} zeigt die Testergebnisse der fünf untersuchten Transformer-Modelle nach dem Fine-Tuning auf dem 
Klassifikationsdatensatz. 
Dargestellt sind Accuracy, F1-Score, Loss sowie die Evaluationszeit des Testdatensatzes. Das Modell XLM-RoBERTa-Large erzielt mit einer 
Accuracy und einem F1-Score von jeweils 97,95\% die besten Resultate, gefolgt von RoBERTa-Large und RoBERTa-Base. 
Das schwächste Ergebnis liefert BERT-Base-Uncased mit einem Accuracy-Wert von 95,33\%. Insgesamt zeigt sich, dass die \textit{Large}-Varianten
der Tranformer Modelle bessere Ergebnisse erzielen, allerdings auch mit einem erhöhten Evaluationsaufwand verbunden sind.

\begin{table}[!ht]
\centering
\begin{tabular}{lcccc}
    \toprule
    \textbf{Modell} & \textbf{Accuracy} & \textbf{F1-Score} & \textbf{Loss} & \textbf{Eval-Zeit (s)} \\
    \midrule
    XLM-RoBERTa-Large & 0.9795 & 0.9795 & 0.1070 & 13.96 \\
    RoBERTa-Large     & 0.9765 & 0.9764 & 0.1395 & 14.04 \\
    RoBERTa-Base      & 0.9751 & 0.9751 & 0.1273 & 5.86 \\
    XLM-RoBERTa-Base  & 0.9717 & 0.9716 & 0.1527 & 5.77 \\
    BERT-Base-Uncased & 0.9533 & 0.9531 & 0.1952 & 6.33 \\
    \bottomrule
\end{tabular}
\caption{Testergebnisse der Transformer Modelle nach dem Fine-Tuning}
\label{tab:vergleich_der_transformer_modelle}
\end{table}

In Tabelle \ref{tab:vergleich_lightgbm_modelle} dargestellt, sind die Metriken verschiedener LightGBM-Modells, welche auf den neu erzeugten 
Embeddings der verschiedenen Transformer-Modellen trainiert wurden.
Dabei wurden jeweils die letzten vier Schichten der Modelle gemittelt und als Input für den LightGBM-Classifier verwendet. 
Auch in diesem Setup erzielt XLM-RoBERTa-Large mit einer Accuracy von 97,84\% und einem F1-Score von 97,73\% die besten Ergebnisse. 
Es folgen die \textit{Base}-Varianten von XLM-RoBERTa und RoBERTa mit sehr ähnlichen Werten. Deutlich schwächer schneidet BERT-Base-Uncased ab, 
was auf eine geringere Qualität der erzeugten Embeddings hinweist. Insgesamt zeigt sich, dass die Qualität der Repräsentationen der Transformer-Modelle 
einen wesentlichen Einfluss auf die nachgelagerte Klassifikationsleistung hat.

\begin{table}[!ht]
\centering
\begin{tabular}{lcc}
    \toprule
    \textbf{Embedding-Modell} & \textbf{Accuracy} & \textbf{F1-Score} \\
    \midrule
    XLM-RoBERTa-Large & 0.9784 & 0.9773 \\
    XLM-RoBERTa-Base  & 0.9753 & 0.9738 \\
    RoBERTa-Large     & 0.9753 & 0.9738 \\
    RoBERTa-Base      & 0.9729 & 0.9713 \\
    BERT-Base-Uncased & 0.9472 & 0.9439 \\
    \bottomrule
\end{tabular}
\caption{Testergebnisse der LightGBM-Modelle nach dem Training auf den Embeddings}
\label{tab:vergleich_lightgbm_modelle}
\end{table}

Tabelle \ref{tab:transformer_vs_lightgbm} vergleicht die Accuracy- und F1-Werte der Transformer und LightGBM Modelle.
Der Differenzwert ($\Delta$) zeigt den Unterschied zwischen den beiden Modellen.
Es zeigt sich, dass die direkt feinjustierten Transformer-Modelle in fast allen Fällen leicht bessere Ergebnisse erzielen als die Kombination aus 
Embeddings und LightGBM.

\begin{table}[!ht]
\centering
\begin{tabular}{lcccccc}
    \toprule
    \multirow{2}{*}{\textbf{Modell}} &
    \multicolumn{2}{c}{\textbf{Transformer}} &
    \multicolumn{2}{c}{\textbf{LightGBM}} &
    \multicolumn{2}{c}{$\Delta$ (LGBM - TF)} \\
    \cmidrule(lr){2-3} \cmidrule(lr){4-5} \cmidrule(l){6-7}
    & \textbf{Accuracy} & \textbf{F1} & \textbf{Accuracy} & \textbf{F1} & $\Delta$ Acc & $\Delta$ F1 \\
    \midrule
    XLM-RoBERTa-Large & 0.9795 & 0.9795 & 0.9784 & 0.9773 & $-$0.0011 & $-$0.0022 \\
    RoBERTa-Large     & 0.9765 & 0.9764 & 0.9753 & 0.9738 & $-$0.0012 & $-$0.0026 \\
    RoBERTa-Base      & 0.9751 & 0.9751 & 0.9729 & 0.9713 & $-$0.0022 & $-$0.0038 \\
    XLM-RoBERTa-Base  & 0.9717 & 0.9716 & 0.9753 & 0.9738 & +0.0036 & +0.0022 \\
    BERT-Base-Uncased & 0.9533 & 0.9531 & 0.9472 & 0.9439 & $-$0.0061 & $-$0.0092 \\
    \bottomrule
\end{tabular}
\caption{Vergleich von Transformer- und LightGBM-Modellen: Accuracy, F1 und Differenz}
\label{tab:transformer_vs_lightgbm}
\end{table}

Ein möglicher Grund für die leicht schlechteren Metriken bei der Verwendung von LightGBM könnte in der hohen semantischen Dichte und Dimension der 
Transformer-Embeddings liegen. Diese enthalten sehr viele komplexe, kontextuelle Informationen, die zwar für überwachte Modelle direkt nutzbar sind, 
klassische Modelle wie LightGBM jedoch nur eingeschränkt verarbeiten können.
\cite{reimers-gurevych-2019-sentence} erklärt, dass die Architektur von BERT für die Anwendung in weiteren unüberwachte Verfahren problematisch ist.
Es lässt sich vermuten, dass Gleiches auch für überwachte Modelle gilt dass Gleiches auch für überwachte Modelle gilt, 
sofern diese nicht speziell auf die Strukturen der Transformer-Embeddings ausgelegt sind.

\section{Vergleich mit verwandten Arbeiten}

Tabelle \ref{tab:vergleich_literatur} vergleicht die Accuracy- und F1-Scores der in dieser Arbeit entwickelten Modelle mit den besten Resultaten zweier 
verwandter Studien. Während \cite{Essa:2023aa} ihre Modelle auf drei unterschiedlichen Fake-News-Datensätzen getestet hat 
(ISOT, TI-CNN, FNC), stammen die Ergebnisse aus \cite{V_G_2024} vom FNDD-Datensatz. In dieser Arbeit wurde ein weiterer selbst entwickelter Datensatz verwendet. 
Aufgrund dieser Unterschiede ist ein direkter Vergleich daher nur bedingt sinnvoll.

Auffällig ist, dass in den Arbeiten von \cite{Essa:2023aa} und \cite{V_G_2024} ausschließlich Ergebnisse für klassifikatorgestützte Ansätze 
(z.,B. BERT-Embeddings mit LightGBM oder LSTM) berichtet werden. Es fehlen dort jedoch vergleichbare Kennzahlen für direkt feinjustierte Transformer-Modelle 
wie sie in dieser Arbeit durchgeführt wurden. Ein direkter Leistungsvergleich mit reinem Transformer-Fine-Tuning ist folglich nicht möglich.

Trotzdem lässt sich beobachten, dass das in dieser Arbeit durchgeführte Fine-Tuning des XLM-RoBERTa-Large-Modells mit einer Accuracy und einem F1-Score von 
jeweils 97,95\% sehr gute Ergebnisse erzielt und damit nahe an den besten LightGBM-Kombinationen der Vergleichsarbeiten liegt. 
Auch das LightGBM-Modell auf Basis der XLM-RoBERTa-Embeddings erreicht mit 97,84\% (Accuracy) und 97,73\% (F1) ein konkurrenzfähiges Niveau.

Die teils noch höheren Werte in \cite{Essa:2023aa} könnten unter anderem auf Unterschiede in den Datensätzen zurückzuführen sein. 
So ist beispielsweise der ISOT-Datensatz stark strukturiert und enthält viele sich wiederholende Phrasen, was Klassifizierungen durch Embedding-Modelle begünstigen kann \cite{s22186970}. 

\begin{table}[ht]
\centering
\begin{tabular}{p{2.4cm} p{3.1cm} p {4cm} cc}
    \toprule
    \textbf{Arbeit} & \textbf{Datensatz} & \textbf{Bestes Modell} & \textbf{Acc. (\%)} & \textbf{F1 (\%)} \\
    \midrule
    \cite{Essa:2023aa}  & ISOT    & BERT + LightGBM            & 99.88 & 99.88 \\
    \cite{Essa:2023aa}  & FNC     & BERT + LightGBM             & 99.06 & 99.05 \\
    \textbf{Diese Arbeit} & Eigener Datensatz & XLM-RoBERTa-Large & 97.95 & 97.95 \\
    \textbf{Diese Arbeit} & Eigener Datensatz & XLM-RoBERTa-Large + LightGBM & 97.84 & 97.73 \\
    \textbf{Diese Arbeit} & Eigener Datensatz & RoBERTa-Large + LightGBM & 97.53 & 97.38 \\
    \cite{Essa:2023aa}  & TI-CNN  & BERT + LightGBM           & 96.94 & 97.42 \\
    \cite{V_G_2024}     & FNDD    & RoBERTa + LightGBM           & 96.54 & 97.60 \\
    \textbf{Diese Arbeit} & Eigener Datensatz & BERT + LightGBM & 94.72 & 94.39 \\
    \bottomrule
\end{tabular}
\caption{Vergleich der erzielten Accuracy- und F1-Scores mit verwandten Arbeiten}
\label{tab:vergleich_literatur}
\end{table}

Zusammenfassend zeigen die Ergebnisse, dass die in dieser Arbeit verfolgten Methoden sowohl im direkten Fine-Tuning als auch im kombinierten
LightGBM-Ansatz mit aktuellen Forschungsarbeiten konkurrenzfähig sind.
Zudem wird aber auch deutlich, dass ein sinnvoller Vergleich verschiedener Modellarchitekturen sehr schwer fällt.
So liefert zum Beispiel die Kombination von BERT und LightGBM in dieser Arbeit mit 94,72\% (Accuracy) das schlechteste Ergebniss, während die gleiche Kombination
in \cite{Essa:2023aa} 5,16\% mehr erzielt.
Mögliche Ursachen hierfür können außerdem die unterschiedlichen Konfigurationen der Hyperparameter beim Modelltraining oder die Bildung der Embeddings aus den
\textit{hidden states} sein. 

\section{Ergebnisse der finalen Anwendung}

Die bereitgestellte Anwendung liefert einen automatischen Klassifizierungsservice für Nachrichtenartikel auf den Seiten 
\textit{BILD.de}, \textit{taz.de} und \textit{spiegel.de}.

Die in Abbildung \ref{fig:screenshot_anwendung} dargestellten Beispiele verdeutlichen, dass die vorgenommenen Klassifizierungen noch Optimierungspotenzial aufweisen.
Ein möglicher Grund hierfür liegt in der unzureichenden Abdeckung aktueller Themen in den für das Training verwendeten Datensätzen. Infolgedessen verfügen die Modelle nicht 
über ausreichenden Kontext, um die abgefragten Domänen adäquat bewerten zu können.

Ein exemplarischer Fall ist der Artikel auf BILD.de mit dem Titel „Trump geht plötzlich auf Israel los“. Obwohl der Inhalt des Artikels faktisch korrekt ist, 
wird er vom Modell fälschlicherweise als falsch eingestuft.

Diese Fehleinschätzung könnte durch semantische oder stilistische Merkmale in der Überschrift bedingt sein. Um diesen Einfluss näher zu untersuchen, wäre es denkbar, 
Titelvariationen wie „Trump kritisiert Israel - eine Kehrtwende in seiner bisherigen Haltung“ zu testen.

\begin{figure}[htbp]
    \begin{center}
        \includegraphics[width=\linewidth]{static/screenshots_anwendung.png}
        \caption{\label{fig:screenshot_anwendung} Screenshots der Anwendung auf den verschiedenen Nachrichtenportalen - faktisch sind alle Artikel korrekt}
    \end{center}
\end{figure}